\BourbakiExercises{习题}

\begin{enumerate}
\def\labelenumi{\arabic{enumi})}
\item
  证明左阿廷且右诺特的环 A 也是右阿廷的（把右 A-模 \(\Re(A)^{k}/\Re(A)^{k+1}\) 视为 \((A/\Re(A))\)-模）。
\item
  设 A 是一个（左）阿廷环，r 是它的根。
\end{enumerate}

\begin{enumerate}
\def\labelenumi{\alph{enumi})}
\item
  证明 A 的每个右理想 b 都含有一个极小右理想（考虑诸理想 \(b \cap r^{p}\)，并注意若 br = 0，则 b 是一个右 A/r-模）。
\item
  证明 A 的左（相应地，右）底座是 t 的右（相应地，左）零化子。由此推出 A 的每个极小双边理想都含于 A 的左、右底座的交中。
\end{enumerate}

\begin{enumerate}
\def\labelenumi{\arabic{enumi})}
\setcounter{enumi}{2}
\tightlist
\item
  设 A 是一个阿廷环，使得环 A/ \(\Re(A)\) 是单的。证明 A 同构于阿廷局部环 B 上的一个矩阵环 \(\mathbf{M}_{r}(\mathrm{B})\)（利用 VIII, 第 169 页的习题 19 与 20，以及 VIII, 第 94 页的习题 19）。
\end{enumerate}

¶ 4) 若对 A 的每个元素 a，存在一个整数 n 与 A 的一个元素 x，使得 \(a^{n}xa^{n}=a^{n}\)，则称环 A 是伪正则的。绝对平坦环（VIII, 第 171 页，习题 27）是伪正则的。

\begin{enumerate}
\def\labelenumi{\alph{enumi})}
\item
  证明伪正则环是一个佐恩环（VIII, 第 171 页，习题 26），其中每个不可逆元素都是左、右零因子。
\item
  证明伪正则环的每个商环都是伪正则的（参见 VIII, 第 171 页，习题 27, d)）。
\item
  证明伪正则环 A 的中心 Z 是伪正则的（证明若 \(a \in Z\) 与 \(x \in A\) 满足 \(a^{n}xa^{n} = a^{n}\)，则 \(a^{2n}x^{k}\) 对每个 \(k \geqslant 1\) 都属于 Z）。
\item
  设 p 是一个素数，\(U_{p}\) 是群 Q/Z 的 p-准素分量，A 是群 \(U_{p}\) 在体 \(F_{p}\) 上的（交换）代数。证明环 A 是伪正则的，但 \(A^{N}\) 不是。
\end{enumerate}

¶ 5) 设 n 是一个整数。若对 A 的每个元素 a，存在一个 \(x \in A\) 使得 \(xa^{n+1} = a^n\)，则称环 A 是 n-正则的。

\begin{enumerate}
\def\labelenumi{\alph{enumi})}
\item
  假设 \(\mathrm{A}\) 是 \(n\)-正则的。证明我们有 \(r^n = 0\)，对每个元素 \(r\)（\(\Re(\mathrm{A})\) 的）皆然。
\item
  证明 n-正则本原环同构于体上的一个矩阵环 \(\mathbf{M}_{r}(\mathrm{D})\)，其中 \(r \leqslant n\)。（注意若 V 是一个维数 \textgreater{} n 的 D-向量空间且 e 是 V 中一个非零向量，则 \(\operatorname{End}_{\mathrm{D}}(\mathrm{V})\) 的每个稠密子环都含有一个元素 u，使得 \(u^{n}(e) \neq 0\) 且 \(u^{n+1}(e) = 0\)。）
\item
  证明 n-正则环是一个佐恩环（VIII, 第 171 页，习题 26）。（首先注意关系式 \(xa^{n+1}=a^{n}\) 蕴含 \(x^{n}a^{2n}=a^{n}\)。利用 b)、VIII, 第 168 页的习题 12, a) 及 VIII, 第 27 页的命题 2，推出我们有 \(a^{n}x^{n}a^{n}-a^{n}\in\Re(A)\)。最后，利用 VIII, 第 159 页的引理 1，证明若 A 的两个元素 y 与 z 满足 \(y=zy^{2}\) 且 \(yzy-y\in\Re(A)\)，则存在一个 \(t\in A\)，使得 \(y=ty^{2}\) 且 \((ty)^{2}=ty.\) 若环 A 此外还是无根的，则它是伪正则的（习题 4）。
\end{enumerate}

\(d\) ) 设 A 是一个除 0 外不含幂零元素的环。则 A 是 \(n\)-正则的当且仅当它是绝对平坦的。

（若 A 是绝对平坦的，利用 VIII, 第 171 页的习题 27, b)。若 A 是 n-正则的，首先按上面的记号注意 \(e = x^{n}a^{n}\) 是幂等的，因而是中心的（VIII, 第 172 页的习题 30, a)），且我们有 ea = a，从而 A 是 1-正则的。然后利用问题 a) 与 b) 证明 A 同构于体的积的一个子环，并推出关系式 \(xa^{n+1} = a^{n}\) 蕴含 axa = a。）

¶ 6) 设 A 是一个阿廷环，\(n = \text{long}_{\text{A}}(\text{A}_s)\)。

\begin{enumerate}
\def\labelenumi{\alph{enumi})}
\item
  证明环 A 是 \(n\)-正则的（习题 5）。
\item
  证明环 A 是伪正则的。（设 r 是使 \(\Re(A)^{r}=0\) 的一个整数，m 是环 B = A/ \(\Re(A)\) 的长度。通过化归到环 B 是单的情形，证明对每个 \(b\in B\)，存在 B 的一个与 b 交换且满足 \(b^{m+1}y=b^{m}\) 的元素 y。由此推出，对每个 \(a\in A\)，存在一个 \(x\in A\)，使得 \(a^{mr}xa^{mr}=a^{mr}\)。）
\item
  设 \(\mathfrak{I}\) 是 A 的非幂零左理想所成之集。证明 \(\mathfrak{I}\) 的极小元素是由不可分解幂等元生成的诸理想 \(Ae\)（利用 \(a\) 及习题 5, \(c\)）。证明 A 的每个左理想都是有限多个这种类型的理想与一个含于根中的左理想的直和。
\item
  A 的左理想 a 由一个幂等元生成当且仅当它是左理想的一个有限族 \((Ae_{i})_{i\in I}\) 的直和，其中诸 \(e_{i}\) 是不可分解的正交幂等元。
\end{enumerate}

¶ 7) 设 A 是一个阿廷环，r 是它的根，\(\overline{A}\) 是环 A/r。设 l 是一个非幂零的不可分解左理想；我们有 l = Ae，其中 e 是 A 的一个幂等元。

\begin{enumerate}
\def\labelenumi{\alph{enumi})}
\item
  证明理想 \(\mathfrak{n} = \mathfrak{l} \cap \mathfrak{r}\) 是 \(\mathfrak{l}\) 的唯一的极大子模。A-模 \(\mathfrak{l} / \mathfrak{n}\) 是单的，且同构于 \(\overline{\mathrm{A}}\overline{e}\)，其中 \(\overline{e}\) 是 \(e\) 在 \(\overline{\mathrm{A}}\) 中的像。
\item
  设 M 是一个 A-模。则 \((Ae)M\) 是 M 的一个子模，且 \(e\big((Ae)M\big)=eM\) 是一个 eAe-模。若 N 是 M 的一个 A-子模，则 eAe-模 \(e(M/N)\) 与 eM/eN 同构。若 A-模 M 是单的，则或者我们有 eM=0，或者 M 同构于 \(\overline{A}\overline{e}\) 且 eM 是一个单 eAe-模。
\item
  证明 eAe-模 eM 的长度等于 A-模 M 的若尔当–赫尔德列中同构于 \(\overline{A}\overline{e}\) 的商的个数。
\item
  设 M 具有有限长度且具有唯一的极大子模 N，并且 M/N 同构于 \(l/n\)。证明 M 同构于 Ae 的一个商模（注意 eM 中不属于 N 的一个元素 x 生成 M）。
\end{enumerate}

¶ 8) 设 A 是一个阿廷环。A-模 \(A_{s}\) 是不可分解左理想的一个有限族 \((Ae_{i})_{i\in I}\) 的直和，其中诸 \(e_{i}\) 是正交幂等元（习题 6, d)）。

\begin{enumerate}
\def\labelenumi{\alph{enumi})}
\tightlist
\item
  若两个有限长度的 A-模 M 与 N 之一的一个若尔当–赫尔德列的商同构于另一个的一个若尔当–赫尔德列的商，则称它们是相连的。设 R 是诸理想 \(Ae_{i}\) 所成之集上使两个相连的理想等价的最细等价关系。R 的等价类 \(\mathcal{B}_{l}\)（ \(l \in L\)）称为块；我们将块 \(B_{l}\) 中诸理想之和记作 \(b_{l}\)。证明诸理想 \(b_{l}\) 是 A 的仅有的不可分解双边理想。
\end{enumerate}

（为看出 \(b_{l}\) 是双边的，注意我们有 \(Ae_{i}xe_{j} \subset Ae_{j}\)，且由习题 7, c)，关系 \(e_{i}Ae_{j} \neq 0\) 蕴含 \(Ae_{i} \equiv Ae_{j} \pmod{R}\)。另一方面，设 \((\mathfrak{b}_{m}^{\prime})_{m \in \mathrm{M}}\) 是以 A 为直和的双边理想的一个有限族。证明每个 \(b_{m}^{\prime}\) 都是诸理想 \(Ae_{i}\) 的一个和，然后证明两个相连的理想 \(Ae_{r}\) 与 \(Ae_{s}\) 属于同一个理想 \(b_{m}^{\prime}\)，为此注意：由习题 7，存在一个指标 h 使 \(e_{h}Ae_{r}\) 与 \(e_{h}Ae_{s}\) 都 \(\neq 0\)。）

\begin{enumerate}
\def\labelenumi{\alph{enumi})}
\setcounter{enumi}{1}
\item
  证明诸右理想 \(e_{i}A\) 是不可分解的（考察它们在 A/ \(\Re(A)\) 中的像）。
\item
  设 A 是交换域 K 上的一个有限次数代数。我们用 \(c_{ij}\)（相应地，\(d_{ij}\)）表示 \(Ae_{j}\)（相应地，\(e_{j}A\)）的若尔当–赫尔德列中同构于 \(\overline{A}\overline{e}_{i}\)（相应地，\(\overline{e}_{i}\overline{A}\)）的商的个数，用 \(r_{i}\) 表示体 \(e_{i}Ae_{i}/e_{i}\Re(A)e_{i}\) 在 K 上的次数。证明等式 \(c_{ij}r_{j}=d_{ij}r_{i}\)（确定 \(e_{i}Ae_{j}\) 在 K 上的维数，并利用习题 7, c)）。
\end{enumerate}

\begin{enumerate}
\def\labelenumi{\arabic{enumi})}
\setcounter{enumi}{8}
\tightlist
\item
  设 A 是一个根 \(\mathfrak{r}\) 含于中心的阿廷环。
\end{enumerate}

\begin{enumerate}
\def\labelenumi{\alph{enumi})}
\item
  证明 A/r 中的每个中心幂等元都可提升为 A 中的中心幂等元（若 e 是 A 的一个幂等元且其类在 A/r 中是中心的，把 e 与 xe - ex 交换这一事实对每个 \(x \in A\) 表达出来）。
\item
  若 \(\mathfrak{r}\) 非零，证明 \(A / \mathfrak{r}\) 是交换的（注意我们有关系式 \((xy - yx)z = 0\)，其中 \(x\in A\)、\(y\in A\)、\(z\in \mathfrak{r}\)）。
\item
  由 a) 与 b) 推出 A 同构于一个半单环与有限多个局部环 \(A_{i}\) 的积，其中体 \(\mathrm{A}_{i}/\Re(\mathrm{A}_{i})\) 是交换的。d) 证明每个单 A-模同构于 A 的一个极小左理想。
\end{enumerate}

¶ 10) 对一个环 A，将 A 的元素全为幂零的伪子环所成之集记作 \(\mathcal{N}(\mathrm{A})\)，将 \(\mathcal{N}(\mathrm{A})\) 的极大元素所成之集记作 \(\mathcal{N}_{\max}(\mathrm{A})\)。

\begin{enumerate}
\def\labelenumi{\alph{enumi})}
\tightlist
\item
  设 D 是一个体，n 是一个整数。证明 \(\mathcal{N}_{\max}(\mathbf{M}_{n}(\mathrm{D}))\) 由诸伪环 \(gTg^{-1}\) 组成，其中 g 是 \(\mathbf{M}_{n}(\mathrm{D})\) 的一个可逆元，T 是由矩阵 \((a_{ij})\in\mathbf{M}_{n}(\mathrm{D})\) 中满足 \(a_{ij}=0\)（对 \(i\geqslant j\)）者组成的集（利用 VIII, 第 14 页的习题 2）。
\item
  设 A 是一个阿廷环，\(\pi:A\to A/\Re(A)\) 是典范映射。证明映射 \(P\mapsto\pi(P)\) 定义了从 \(\mathcal{N}_{\max}(A)\) 到 \(\mathcal{N}_{\max}(A/\Re(A))\) 的一个双射。由 a) 推出 \(\mathcal{N}_{\max}(A)\) 的诸元素之交等于 \(\Re(A)\)
\end{enumerate}

并且 \(\mathcal{N}_{\mathrm{max}}(\mathrm{A})\) 的任意两个元素都可通过 A 的一个内自同构互相变换。

\begin{enumerate}
\def\labelenumi{\arabic{enumi})}
\setcounter{enumi}{10}
\tightlist
\item
  设 A 是一个交换环，B 是一个 A-代数，M 是一个阿廷 B-模。假设模 M 是有限生成的。
\end{enumerate}

\begin{enumerate}
\def\labelenumi{\alph{enumi})}
\item
  证明环 \(B_{M}\) 是左阿廷的（利用 VIII, 第 9 页的命题 6），然后证明 B-模 M 具有有限长度。
\item
  证明若 M 是一个单 B-模，则 A-模 M 是同型的（由 VIII, 第 81 页的引理 1 推出同态 \(A \rightarrow \text{End}_{B}(M)\) 的核是一个极大理想）。
\item
  由此得出 M 是一个有限长度的 A-模。
\end{enumerate}

